%=====================================================================
\chapter{Data Science for the Internet of Things}\label{ch:iot}
%=====================================================================
\locator{5}{Part V --- Data Science in Systems Context}

\begin{outcomes}
By the end of this chapter you will be able to:
\begin{tight}
  \item \textbf{Describe} the IoT reference architecture and where analytics sits within it.
  \item \textbf{Characterise} telemetry data and the failure modes specific to sensors.
  \item \textbf{Design} a predictive maintenance system end to end, from framing to deployment.
  \item \textbf{Apply} sensor fusion and explain why multiple weak sensors beat one good one.
  \item \textbf{Evaluate} an IoT model against operational rather than statistical criteria.
\end{tight}
\end{outcomes}

\begin{prereq}
The four chapters this one is built on: \chref{timeseries} (streaming, windowing,
drift), \chref{unsupervised} (anomaly detection without labels), \chref{features}
(window aggregations) and \chref{cloudedge} (edge inference). This is the chapter
those were preparing you for.
\end{prereq}

\begin{keyterms}
telemetry $\bullet$ sensor $\bullet$ actuator $\bullet$ MQTT $\bullet$ gateway
$\bullet$ device twin $\bullet$ sampling rate $\bullet$ sensor drift $\bullet$
calibration $\bullet$ sensor fusion $\bullet$ predictive maintenance $\bullet$
remaining useful life $\bullet$ digital twin $\bullet$ over-the-air update
\end{keyterms}

\section{The Architecture, and Where Data Science Lives in It}

\dsfig{%
\begin{tikzpicture}[font=\scriptsize,
  lay/.style={rounded corners=3pt, draw=#1, fill=#1!8, text=#1!50!black,
              text width=2.35cm, align=center, font=\tiny\bfseries, minimum height=1.45cm},
  arr/.style={-{Stealth[length=2mm]}, gray!60, line width=0.9pt},
  ds/.style={font=\tiny, text=accent!75!black, align=center, text width=2.5cm}]

\node[lay=secondary] (dev)  at (0,0)    {\faIcon{microchip}\\DEVICE\\{\tiny sensors,\\actuators}};
\node[lay=tealhl]    (net)  at (2.9,0)  {\faIcon{wifi}\\CONNECTIVITY\\{\tiny MQTT, LoRa,\\5G}};
\node[lay=goldhl]    (gw)   at (5.8,0)  {\faIcon{server}\\GATEWAY\\{\tiny fog tier}};
\node[lay=primary]   (plat) at (8.7,0)  {\faIcon{database}\\PLATFORM\\{\tiny ingest, store,\\device mgmt}};
\node[lay=purplehl]  (app)  at (11.6,0) {\faIcon{chart-line}\\APPLICATION\\{\tiny dashboards,\\work orders}};

\foreach \a/\b in {dev/net, net/gw, gw/plat, plat/app}{\draw[arr] (\a) -- (\b);}
\draw[arr, gray!45, dashed] (plat.south) .. controls +(0,-1.1) and +(0,-1.1) .. (dev.south);
\node[font=\tiny, text=gray!55!black] at (4.35,-1.55) {control \& over-the-air model updates};

% where data science acts
\node[ds] at (0,1.55)    {\textbf{edge inference}\\threshold, tiny model};
\node[ds] at (5.8,1.55)  {\textbf{fog analytics}\\cross-device correlation};
\node[ds] at (8.7,1.55)  {\textbf{training}\\fleet-wide learning};
\node[ds] at (11.6,1.55) {\textbf{decision support}\\RUL, scheduling};
\foreach \x in {0,5.8,8.7,11.6}{
  \draw[-{Stealth[length=1.6mm]}, accent!60, dashed] (\x,1.05) -- (\x,0.78);}

\node[anchor=north west, align=left, text width=13.8cm, font=\scriptsize, text=primary!60!black]
     at (-1.5,-2.55)
     {\textbf{Read this diagram as a constraint list, not a picture.} Each layer imposes a
      limit that shapes the model: the device limits compute and power, connectivity
      limits bandwidth and guarantees only at-least-once delivery, the gateway is where
      cross-device context first exists, and the platform is the only place with enough
      history to train. A model design that ignores any of these will not deploy.};
\end{tikzpicture}}%
{The IoT reference architecture, annotated with where data science actually runs. Analytics
is not a single box at the end --- it is distributed across four tiers with different
capabilities.}{iot-architecture}

\section{What Telemetry Data Is Like}

\rowcolors{2}{white}{lightbg}
\begin{longtable}{@{}p{3.2cm}p{4.6cm}p{6.4cm}@{}}
\toprule
\rowcolor{primary!15}
\textbf{Property} & \textbf{Consequence} & \textbf{What you must do} \\
\midrule
\endhead
Very high volume & Cannot store or query raw indefinitely & Pre-aggregate to minute/hour grains; retain raw only briefly (\chref{dataeng}) \\
Highly autocorrelated & Effective sample size $\lll$ row count & Never trust a $p$-value computed on raw readings (\chref{inference}) \\
Irregular timestamps & Clocks drift across a fleet & Resample to a fixed grid; use the gateway's clock, not the device's \\
Duplicates & At-least-once delivery guarantees them & Idempotent ingestion keyed on (device, timestamp) \\
Out of order & Network delay and buffering & Watermarks; decide how long to wait (\chref{timeseries}) \\
Missing bursts & Connectivity loss & Forward fill \emph{with a cap}, plus a \texttt{was\_missing} flag \\
Extremely imbalanced labels & Equipment rarely fails & Anomaly detection, not classification (\chref{unsupervised}) \\
\bottomrule
\end{longtable}

\begin{alertbox}[The Failure Mode That Catches Everyone]
A sensor that fails often reports a \emph{perfectly plausible constant}. Nothing is
null, nothing is out of range, no error is logged --- and the model sees a machine
in flawless health. Add a check for \emph{variance too low} alongside your range
checks. A temperature that has not changed by 0.01 degrees in six hours is not a
stable process; it is a dead thermistor.
\end{alertbox}

\dsfig{%
\begin{tikzpicture}
\begin{groupplot}[
  group style={group size=2 by 2, horizontal sep=1.3cm, vertical sep=1.0cm},
  width=6.5cm, height=3.3cm,
  axis lines=left, tick label style={font=\tiny}, label style={font=\tiny},
  title style={font=\tiny\bfseries}, xtick=\empty,
]
\nextgroupplot[title={\color{greenhl}Healthy: noisy around a set point}, ymin=68, ymax=76]
\addplot[greenhl, line width=0.7pt, domain=0:100, samples=180]
  {72 + 1.4*sin(deg(x/3.1)) + 0.9*sin(deg(x*1.7))};

\nextgroupplot[title={\color{accent}Stuck sensor: variance collapses}, ymin=68, ymax=76]
\addplot[accent, line width=0.7pt, domain=0:45, samples=90]
  {72 + 1.4*sin(deg(x/3.1)) + 0.9*sin(deg(x*1.7))};
\addplot[accent, line width=1.1pt, domain=45:100, samples=10] {71.6};
\draw[gray!60, dashed] (axis cs:45,68) -- (axis cs:45,76);
\node[font=\tiny, accent!80!black, anchor=west] at (axis cs:48,74.6) {no error raised};

\nextgroupplot[title={\color{goldhl!60!black}Calibration drift: slow baseline shift}, ymin=68, ymax=80]
\addplot[goldhl, line width=0.7pt, domain=0:100, samples=180]
  {72 + 0.055*x + 1.2*sin(deg(x/3.1))};
\draw[gray!55, dashed] (axis cs:0,72) -- (axis cs:100,72);
\node[font=\tiny, goldhl!50!black, anchor=west] at (axis cs:8,78.6) {the process is fine; the sensor is not};

\nextgroupplot[title={\color{purplehl}Genuine degradation: variance grows}, ymin=64, ymax=80]
\addplot[purplehl, line width=0.7pt, domain=0:100, samples=200]
  {72 + (0.6+0.055*x)*sin(deg(x*2.3)) + 0.5*sin(deg(x/2.7))};
\node[font=\tiny, purplehl, anchor=west] at (axis cs:5,78) {this is the one you want to catch};
\end{groupplot}
\end{tikzpicture}}%
{Four telemetry signatures. Distinguishing a failing \emph{sensor} (top right, bottom left)
from a failing \emph{machine} (bottom right) is the central diagnostic skill in IoT
analytics --- and range checks alone cannot do it.}{telemetry-signatures}

\section{Sensor Fusion}

\begin{definitionbox}[Sensor Fusion]
Combining readings from several sensors to produce an estimate better than any one
of them. Sensors fail and drift independently, so their errors partly cancel --- the
same argument that makes ensembles work (\chref{features}), applied to hardware.
\end{definitionbox}

\begin{worked}[Why Three Cheap Sensors Beat One Good One]
A bearing is monitored for overheating.

\smallskip
\textbf{Single sensor.} One precise thermometer, $\sigma = 0.5^\circ$C. If it
sticks (\dsref{telemetry-signatures}, top right), you have no monitoring at all and
no way to know.

\textbf{Three cheap sensors,} each $\sigma = 1.5^\circ$C, averaged. Because their
errors are independent, the standard error of the mean is
\[
\frac{1.5}{\sqrt{3}} = 0.87^\circ\text{C}
\]
--- worse than the single precise sensor on accuracy alone.

\smallskip
\textbf{But now add the disagreement check.} If one sensor reads $71.4$ and the
other two read $77.2$ and $77.6$, the odd one out is identifiable and can be
excluded. With a single sensor that comparison does not exist.

\smallskip
\textbf{And add a different physical quantity.} Combine temperature with vibration
and current draw. A failing bearing raises all three; a failing \emph{thermistor}
moves only temperature. This is the decisive advantage: fusion across
\emph{different physics} distinguishes a broken machine from a broken sensor, which
no amount of precision in one channel can achieve.

\smallskip
\textbf{Conclusion:} the design goal is not the most accurate sensor but the most
\emph{diagnosable} sensor set. Redundancy buys fault detection; diversity buys
attribution.
\end{worked}

\section{Predictive Maintenance End to End}

\dsfig{%
\begin{tikzpicture}[font=\scriptsize,
  s/.style={rounded corners=3pt, draw=#1, fill=#1!8, text=#1!50!black,
            text width=2.5cm, align=center, font=\tiny\bfseries, minimum height=0.95cm},
  arr/.style={-{Stealth[length=2mm]}, primary!55, line width=1pt},
  note/.style={font=\tiny\itshape, text=gray!55!black, align=center, text width=2.7cm}]

\node[s=primary]   (a) at (0,0)     {1. FRAME\\{\tiny \chref{lifecycle}}};
\node[s=secondary] (b) at (3.3,0)   {2. INGEST\\{\tiny \chref{dataeng}}};
\node[s=tealhl]    (c) at (6.6,0)   {3. WINDOW\\FEATURES\\{\tiny \chref{features}}};
\node[s=goldhl]    (d) at (9.9,0)   {4. MODEL};
\node[s=purplehl]  (e) at (13.0,0)  {5. THRESHOLD\\{\tiny \chref{evaluation}}};
\node[s=greenhl]   (f) at (3.3,-2.4){6. DEPLOY\\TO EDGE\\{\tiny \chref{cloudedge}}};
\node[s=accent]    (g) at (6.6,-2.4){7. MONITOR\\DRIFT\\{\tiny \chref{mlops}}};
\node[s=primary]   (h) at (9.9,-2.4){8. RETRAIN};

\foreach \x/\y in {a/b,b/c,c/d,d/e}{\draw[arr] (\x) -- (\y);}
\draw[arr] (e.south) .. controls +(0,-1.0) and +(2.6,0) .. (f.east);
\draw[arr] (f) -- (g);
\draw[arr] (g) -- (h);
\draw[arr] (h.east) .. controls +(1.6,0) and +(0,-1.4) .. (d.south);

\node[note] at (0,-1.15)   {``will this fail in 30 days?'' --- horizon set by parts lead time};
\node[note] at (6.6,-1.15) {rolling mean, sd, max, rate of change, FFT bands};
\node[note] at (9.9,-1.15) {isolation forest if no failure labels; boosting if some};
\node[note] at (13.0,-2.4) {loop closes: the model is never finished};
\end{tikzpicture}}%
{Predictive maintenance as an assembly of earlier chapters. Nothing here is new
technique --- the contribution of this chapter is the ordering, the constraints and the
operational criteria.}{predictive-maintenance}

\begin{worked}[Framing a Predictive Maintenance Model Properly]
Applying the six framing questions of \chref{lifecycle} to 400 motors.

\smallskip
\textbf{1. Decision:} whether to schedule a motor for inspection in this week's
maintenance window. \emph{Not} ``predict failure'' --- there is a specific action
with a specific capacity.

\textbf{2. Unit:} one motor on one day. So the modelling table has
$400 \times 365 = 146{,}000$ rows per year, not 12 billion raw readings.

\textbf{3. Target:} does this motor fail within the next 30 days? The horizon comes
from parts lead time (3 weeks) plus scheduling slack, not from the data.

\textbf{4. Timing:} only telemetry up to and including today. Any feature computed
over the failure window is a leak (\chref{wrangling}) --- and in this domain the
leak is easy to introduce, because maintenance records are timestamped when they
are \emph{entered}, often after the failure.

\textbf{5. Criterion:} prevent $\geq 80\%$ of unplanned stoppages with $\leq 5$
unnecessary teardowns per 100 motors per year. Note this is recall-at-fixed-precision,
expressed in the operations manager's units.

\textbf{6. Baseline:} fixed-interval servicing every 90 days, which currently
catches about 45\% of failures and performs 4 inspections per motor per year. The
model must beat \emph{that}, not beat random.

\smallskip
\textbf{What this framing prevents.} Without question 2, someone builds a model on
raw readings and reports 99.99\% accuracy (nearly every reading is from a healthy
motor). Without question 6, a model catching 60\% of failures sounds impressive
until you notice the calendar already catches 45\% for free.
\end{worked}

\begin{alertbox}[Evaluate on Operational Criteria, Not Statistical Ones]
Report these four numbers, in this order:
\begin{tightnum}
  \item \textbf{Unplanned stoppages prevented per year} --- the business metric.
  \item \textbf{Unnecessary inspections per year} --- the cost of false positives,
        in technician-days.
  \item \textbf{Warning lead time} --- a correct prediction 2 days ahead is useless
        if parts take 3 weeks. A model with lower recall but longer lead time is
        often the better model.
  \item \textbf{Alerts per technician per week} --- if this exceeds what the team
        can act on, the system will be ignored regardless of its PR-AUC.
\end{tightnum}
Point 3 is the one students miss: for maintenance, \emph{when} you are right
matters as much as \emph{whether}.
\end{alertbox}

\begin{definitionbox}[Digital Twin]
A live software model of a physical asset, continuously updated from its telemetry.
It supports asking counterfactual questions --- ``what happens if we run this pump
20\% faster?'' --- without touching the equipment, and provides a simulated baseline
against which real behaviour can be compared.
\end{definitionbox}

%---------------------------------------------------------------------
\begin{fourlenses}{IoT Analytics}
\lensIOT{This chapter is the IoT lens in full. The distilled version: aggregate
before you model, treat sensor faults as a distinct class from machine faults, place
inference where the latency budget demands, and evaluate on stoppages prevented
rather than on AUC.}

\lensLA{IoT enters the campus as occupancy sensors, lab equipment usage, library
gate counts and environmental monitoring. Room-occupancy telemetry against timetable
data reveals which sessions are genuinely attended --- objectively, without a
register. \textbf{Deploy this carefully:} sensing that can identify individuals'
movements is surveillance, and requires the consent and proportionality analysis of
\chref{ethics}, not merely a privacy notice.}

\lensPM{IoT projects are unusually hard to manage and worth studying as a project
type. They combine hardware lead times, firmware, connectivity contracts and
analytics, so the critical path runs through procurement rather than code. Pilot on
20 devices before committing to 4{,}000 --- the failure modes of a fleet are not
visible at bench scale, and the cost of discovering them after deployment is
measured in truck rolls.}

\lensSEC{IoT devices are a serious and growing attack surface: default credentials,
firmware that is never patched, and plaintext protocols. Two data science
consequences follow. First, \textbf{compromised devices produce anomalous
telemetry}, so your maintenance anomaly detector is also, incidentally, an intrusion
detector --- and you should treat unexplained anomalies as a security question, not
only an engineering one. Second, telemetry can be \emph{forged}: an attacker who
controls a device can send readings that look healthy while the machine is damaged,
or that trigger unnecessary shutdowns. Sensor fusion across independent physics is a
partial defence, because forging three correlated channels consistently is much
harder than forging one.}
\end{fourlenses}

%---------------------------------------------------------------------
\begin{lab}[Telemetry to Alert]
\begin{lstlisting}[language=Python]
import pandas as pd, numpy as np
from sklearn.ensemble import IsolationForest

# --- 1. resample to a fixed grid; device clocks cannot be trusted -------
t = (raw.set_index("gateway_ts")            # gateway clock, not device clock
        .groupby("device_id")
        .resample("1min")
        .agg(temp=("temp", "mean"), vib=("vib", "mean"),
             current=("current", "mean"), n=("temp", "size")))

# --- 2. sensor health checks BEFORE any modelling ----------------------
w = 60  # one hour
health = t.groupby("device_id").rolling(w, min_periods=30).agg(["std", "mean"])

stuck   = health[("temp", "std")] < 0.01        # dead sensor, not a stable process
missing = t["n"] == 0                            # no readings arrived at all
drifted = (health[("temp", "mean")]
           - health[("temp", "mean")].groupby("device_id").transform("first")).abs() > 5

t["sensor_suspect"] = stuck | missing | drifted  # exclude from training; alert separately

# --- 3. window features (Chapter 15) -----------------------------------
g = t.groupby("device_id")
feat = pd.DataFrame(index=t.index)
for col in ("temp", "vib", "current"):
    feat[f"{col}_mean_1h"]  = g[col].transform(lambda s: s.rolling(60).mean())
    feat[f"{col}_std_1h"]   = g[col].transform(lambda s: s.rolling(60).std())
    feat[f"{col}_max_24h"]  = g[col].transform(lambda s: s.rolling(1440).max())
    feat[f"{col}_slope_24h"] = g[col].transform(
        lambda s: s.rolling(1440).apply(lambda v: np.polyfit(range(len(v)), v, 1)[0],
                                        raw=True))
# rising variance is the degradation signature of Figure 25.2, bottom right
feat["vib_var_ratio"] = feat["vib_std_1h"] / g["vib"].transform(
    lambda s: s.rolling(10080).std())          # this hour vs this week

# --- 4. unsupervised: there are no failure labels ----------------------
train = feat[~t["sensor_suspect"]].dropna()
iso = IsolationForest(contamination=0.005, random_state=0).fit(train)
feat["score"] = -iso.score_samples(feat.fillna(0))

# --- 5. alert on PERSISTENCE, not on a single reading ------------------
# One anomalous minute is noise. Six of the last ten is a signal.
flag = feat["score"] > np.percentile(feat["score"], 99.5)
feat["alert"] = flag.groupby(level="device_id").transform(
    lambda s: s.rolling(10).sum() >= 6)

print(feat.groupby(level="device_id")["alert"].sum().sort_values(ascending=False).head())
\end{lstlisting}

\textbf{Step 5 is the one that makes this deployable.} Alerting on every anomalous
reading generates hundreds of alerts a day and the system is switched off within a
week. Requiring persistence trades a little detection latency for a false-alarm rate
technicians will tolerate --- the practical form of the base-rate argument from
\chref{probability}.
\end{lab}

%---------------------------------------------------------------------
\begin{checkpoint}
\begin{tightnum}
  \item A temperature sensor has reported 71.6$^\circ$C exactly for eight hours. Range
        and null checks both pass. What is wrong, and what check catches it?
  \item Why is ``one row per raw reading'' almost never the right unit of analysis
        for a predictive maintenance model?
  \item A model predicts failures with 85\% recall but only 24 hours ahead, while
        parts take three weeks to arrive. Is it useful?
\end{tightnum}
\end{checkpoint}

%---------------------------------------------------------------------
\begin{chaptersummary}
\begin{tight}
  \item IoT analytics is distributed across device, gateway, platform and
        application, each with different compute, latency and context.
  \item Telemetry is high-volume, autocorrelated, duplicated, out of order and
        almost never labelled --- so aggregate first and expect to work unsupervised.
  \item Distinguishing a failing sensor from a failing machine is the core
        diagnostic skill; a stuck sensor passes every conventional validity check.
  \item Sensor fusion across \emph{different physical quantities} is what makes that
        distinction possible.
  \item Frame predictive maintenance around the maintenance decision, at a device-day
        grain, with the horizon set by parts lead time and the baseline set by the
        existing service interval.
  \item Evaluate on stoppages prevented, unnecessary inspections, warning lead time
        and alerts per technician --- and alert on persistence, not on single readings.
\end{tight}
\end{chaptersummary}

\begin{reviewq}
  \item \textbf{(MCQ)} At-least-once delivery in an IoT platform means ingestion must
        be: (a)~fast (b)~idempotent (c)~encrypted (d)~batched.
  \item \textbf{(MCQ)} Which signature indicates genuine mechanical degradation
        rather than sensor failure? (a)~variance collapsing to zero (b)~a slow
        baseline shift in one channel (c)~growing variance across correlated channels
        (d)~missing readings.
  \item \textbf{(Short)} Explain why warning lead time can matter more than recall in
        predictive maintenance.
  \item \textbf{(Short)} Give two reasons predictive maintenance is usually an
        unsupervised problem in practice.
  \item \textbf{(Short)} How does sensor fusion across different physical quantities
        help distinguish a broken sensor from a broken machine?
  \item \textbf{(Applied)} Design a predictive maintenance system for 150 HVAC units
        in a university. Answer all six framing questions, name the features you would
        build, state where inference runs and why, and give the four numbers you would
        report to the estates manager.
  \item \textbf{(Applied)} An attacker compromises 30 devices and makes them report
        healthy readings while the equipment degrades. Describe three data-side
        defences, and explain which of them your existing anomaly detector already
        provides.
\end{reviewq}
